\chapter{知识迁移与共享的评价}
\label{chap:ket-evaluation}

学完迁移和共享的各种方法，最后必须面对一个毫不浪漫的问题：它们真的有用吗？而“有没有用”在多任务、迁移、融合这些场景里，几乎从来不是一个总分能回答的。先看三个特别容易被漂亮数字掩盖的陷阱。其一，一个迁移到新植物园的模型，在只有少量标注时准确率明显更高，可一旦给足数据充分训练，反而可能不如从零训练——早期的胜利不等于全程领先。其二，物种和病害一起训练后平均成绩上升了，但这个平均可能是用“病害漏检增多”换来的，平均变好、关键任务却变差。其三，两个现成模型各自都能正确回答老问题，也绝不代表把它们组合起来就能回答一个全新的联合问题，比如“找出指定物种中患病的那些花”——分别会，不代表合起来会。因此评价必须拆开三件事分别检查：迁移给目标学习带来了多少收益、多任务之间谁受益谁受损、以及融合是否真的产生了超出简单拼接的新组合能力。

先划清本章和前面评价章的分工：第~\ref{chap:ket-evolution-evaluation}章处理的是时间维度上的问题——跨阶段能力保持、后续学习、事实修正、训练影响的消除；本章承接第~\ref{chap:ket-transfer-adaptation}章的迁移和第~\ref{chap:ket-sharing}章的共享，集中评价“新任务迁移”和“多任务共享”，并顺带检查第~\ref{chap:knowledge-fusion}章的多来源融合是否既保住了各来源的原能力、又支持了联合请求。全章遵循一条严谨的评价链条：先把指标和主张定清楚（15.1），再安排能让对应条件真正出现的评价样本（15.2），然后设计访问权限、直接对照和重复协议（15.3），最后才分析收益分布、资源代价和结论能推广到哪里（15.4）。

为了让“好”和“坏”方向统一，本章沿用式~\eqref{eq:ket-reuse-risk-gain}的约定：$R_T$ 是目标分布上的期望损失，$f_0$ 是不使用待检验来源的对照方案，$f_K$ 是使用该来源的方法，定义收益 $\Delta_T=R_T(f_0)-R_T(f_K)$——对照损失减去方法损失，所以收益为正代表方法确实降低了损失。遇到准确率、召回率这类“越大越好”的增大型指标，则反过来用方法分数减对照分数，同样保证“正即改善”。方向统一后才能把不同指标的收益并列讨论。最后声明：本章所有花卉照片上的成绩、任务数字和费用都是用于复算定义的教学构造，不代表任何实测性能；初始学习曲线和来源准备费的设定取自第~\ref{chap:ket-reuse-selection}章，而引用的基准研究只用来说明真实协议长什么样、结论边界在哪里。

\section{评价目标与指标的确定}
\label{sec:ket-evaluation-metrics}

评价的第一步不是跑分，而是把“主张”写成可被数据判定对错的句子，因为不同主张对应完全不同的实验。设想两类典型要求：一类是“在只允许使用 80 个适应标签的前提下，把新地区错误率压到最低”——预算被钉死，比的是这个预算点上的效果；另一类是“必须达到 90% 准确率”——效果被钉死，比的是各方案为达标各自花了多少标签。这是两项各自独立、可分别被结果否定的主张，不能用同一组数字含糊带过。多任务也一样：“平均成绩改善”绝不等于“允许关键的病害任务退步”，它们是不同的承诺。所以每一项主张都应明确写清四件事：评价的对象是什么、对照是谁、在什么预算下、达标条件是什么。多个目标可以按事先约定的权重汇总成一个总分，但各项的原始结果必须原样保留——否则一旦某个关键任务被其他任务的改善在平均里抵消，你将无从察觉。

\subsection{目标表现与适应效率的度量}
\label{sec:ket-evaluation-target}

先处理“固定预算比效果”这类主张。做法是让两个方案在同一组事先预定的预算点上接受测量，而不是各自挑自己最好看的时刻。用 $s_r(b)$ 记录第 $r$ 次独立运行在预算 $b$ 处得到的目标分数；横轴的预算可以是标签数、更新步数或设备时间，但一旦选定就必须全程口径一致。这里要特别防止一种作弊式比较：方案 A 报它在 80 标签的成绩，方案 B 却报它在 200 标签的成绩，然后相减——这比的根本不是同一件事。为保证可复查，每次运行都应保存样本划分、随机种子、来源模型版本、实际可见的目标样本与描述、分数按什么计数，以及累计消耗的资源。

为了公平且充分利用数据，最好采用配对设计：让两个方案在完全相同的目标样本和可匹配的随机条件下各跑一次，然后先在“同一对运行”内部求收益 $d_r(b)$（增大型指标取方法减对照，损失类指标取对照减方法，都保证正为好）。配对的好处是能消掉大量与方案无关的共同波动——这一批测试图特别难，会同时压低两个方案的成绩，但它们的差值受影响小得多，于是配对差往往比两组独立成绩更能显出真实差异。要同时描述“平均收益有多大”和“这个收益在不同运行间稳不稳”，对 $R$ 个配对差求均值和样本标准差：
\begin{equation}
  \bar d(b)=\frac1R\sum_{r=1}^{R}d_r(b),\qquad
  s_d(b)=\sqrt{\frac{1}{R-1}
    \sum_{r=1}^{R}\bigl(d_r(b)-\bar d(b)\bigr)^2}.
  \label{eq:ket-evaluation-run-summary}
\end{equation}
这里要求运行次数 $R>1$（一次运行无从谈波动）。$s_d$ 刻画的是配对收益本身在各次运行间的离散程度，提醒你收益稳不稳；但要注意它只是差值的样本标准差，并不直接等于均值的置信区间（那还要除以 $\sqrt R$，后面 15.3 节专门讲）。正因为配对保留了“同批困难样本同时影响两方案”的联系，如果你反过来从两个方案里各挑一次最好的运行相减，就彻底破坏了同条件对应，得到的差值没有意义。最后一条纪律：用于调参的验证结果和方案锁定后的独立测试结果必须分开，不能用测试集挑出来的成绩冒充泛化能力。

图~\ref{fig:ket-reuse-learning-curves} 给了一个配对测量的完整算例，值得逐点读。在 20、40、80、160 个适应标签四个预算点上，迁移方案准确率为 $(82,88,91,92)\%$，从头训练为 $(65,76,85,91)\%$，两者始终在同一组 100 张验证照片上评测。逐点相减，迁移的优势是 $(17,12,6,1)$ 个百分点——一个非常清楚的故事：迁移在数据极少时帮助最大（17 个点），随着标签增多优势一路收窄，到 160 标签时几乎追平。这正是“早期领先不代表全程领先”的量化体现。还要注意横轴的口径：这里只计了适应标签，没有把来源筛选和验证用掉的标签算进去；若主张是“比总开发标签”，就得按各方案实际使用的全部标签重画横坐标，而不能只给轴改个名字。

单点单跑说明不了稳定性，所以再补两次运行，仍用相同预算点和同一组 100 张照片。第二次迁移为 $(80,87,90,91)\%$、从头为 $(64,75,84,90)\%$；第三次迁移 $(81,86,89,89)\%$、从头 $(66,77,85,89)\%$。聚焦 80 标签这个点，三次的配对收益为 $(6,6,4)$ 个百分点，代入式~\eqref{eq:ket-evaluation-run-summary} 得均值约 $5.33$、样本标准差约 $1.15$ 个百分点——收益集中在 4 到 6 之间，比较稳。两个计数误区点一下：第一次那条曲线在 80 标签的 6 个点优势只是三次运行中的一次，不能拿它代替整体的均值和波动；另外，三次运行总共做了 300 次预测，但这不是 300 张独立照片——还是那 100 张被重复评了三遍，独立样本数没有增加。

现在转到第二类主张——“达标要花多少”。这时要从学习曲线上读出：方案第一次达到预定水平，用掉了多少预算。对增大型分数，先固定达标阈值 $s_\star$ 和一组实际观测过的预算点 $B$，再定义第 $r$ 次运行的“经验到达量”（首次过线的预算）：
\begin{equation}
  b_{\star,r}=\min\{b\in B:s_r(b)\geq s_\star\}.
  \label{eq:ket-evaluation-threshold}
\end{equation}
这个最小值只在确实存在过线点时才有意义。如果跑到最大预算 $b_{\max}$ 都没达标，就应如实记成“截至 $b_{\max}$ 未达到”，绝不能把终点的预算硬填成它的成功成本（那等于假设它在终点成功了）。这种“观察期已经结束、但关心的事件（首次过线）还没被记录到”的情形，在统计学里叫\term{右删失}（right censoring）——就像追踪一种药何时起效，研究结束时某病人还没起效，你只知道“起效时间晚于观察终点”，不知道具体是多久。它既不告诉你未测量的中间预算上成绩如何，也不保证继续给预算就一定能过线。

式~\eqref{eq:ket-evaluation-threshold} 用的是最简单的“单点过线”规则（一个点达标就算）。如果你的上线标准更稳健，比如要求连续几个预算点都达标、或要求置信区间下界都过线，那就必须在开测之前把判据改好，而不能看到数据后再挑对自己有利的规则。还要记住验证集在这里的角色：它只决定“这次运行在哪个点停止/达标”，据此选出的模型仍要拿到独立测试集上评价，经验阈值本身并不保证总体准确率真的达到 $s_\star$。

把前面三次运行按“首次达到 90% 的预算”汇总：迁移方案是 $(80,80,>160)$，从头训练是 $(160,160,>160)$，其中 $>160$ 表示到观测终点仍未记录到过线（右删失）。读这组数字要分层：截至 160 标签，两方案都只有 $2/3$ 的运行成功；但截至 80 标签，迁移已有 $2/3$ 成功，从头训练是 0——这才是迁移在“快速达标”上的真实优势。统计时必须保留那个失败的第三次运行。一个常见错误是“只对成功的运行求平均”：那样会得到迁移平均 80、从头平均 160，看起来干净，实则是只对“本来就会成功”的运行取的条件均值，系统性隐瞒了失败率，绝不能称为“全部运行的平均成本”。

右删失数据在计算上还有一种约定处理：事先规定把未达标者按观察上限 $b_{\max}$ 截断（即姑且记成 160），于是迁移得到 $(80+80+160)/3\approx106.67$、从头训练得到 160。这样能算出数，但命名必须准确——这叫“截断预算统计量”，其中第三个 160 只是“停止观察时已经允许投入的预算”，并不是真测到达标的成本。它既不是无删失情况下的真实平均到达量，也不提供“第三次运行最终会成功”的任何证据。用的时候要把这层含义讲清楚。

到达量的精度受观测间隔限制：第一条曲线只能确认迁移“在 80 标签这个已测点已经过线”，真实的首次过线点可能在 40 到 80 之间，插值补不出精确成本。若想用曲线下面积（AUC）把整条学习曲线汇成一个数，也得先规定横轴是标签数 $n$ 还是 $\log n$、积分区间和插值方法；尤其要注意，在 20、40、80、160 这种等比间隔点上把分数简单平均，并不等于按标签投入做积分（两点之间的横坐标宽度不同，等比点尤其如此）。最后，跨任务汇总还有单位鸿沟：准确率（百分点）、定位误差（像素）、推理成功率（比率）方向和量纲都不同，必须先各自定好收益方向、再按预先约定的尺度和权重比较，不能直接相加。

评价零样本和少样本时，要把“到底用了什么信息”列清楚，否则“零样本”这个词很容易名不副实。零样本预测时类别名称、属性描述、预训练模型当然都可以参与，但如果你事先用目标标签去挑选过提示模板，那这份监督就必须报告——不能因为预测那一刻没有更新模型参数，就声称“完全没有目标监督”：选提示用掉的标签也是目标信息。少样本评价则至少应列明：分几个类、每类几个支持样本、查询集多大，以及是否还看到了额外的无标签目标输入。这些条件不同，“同样是 5-way 1-shot”的难度和结论可能完全不同。

当新物种压根不在模型原来的输出空间里时，一件事被错误地打包成了三件不同的能力，评价必须分开。第一是发现未知：面对没见过的对象，能不能识别出“这不在已知目录里”，这要同时报告两类错误——把已知类误拒成未知的比率、和把未知类误接纳成某个已知类的比率（呼应第13章的拒识阈值）。第二是学习新类：在给了新类样本后，能不能在新类的留出照片上正确分类。第三是校准：把旧类和新类放在一起、让它们在同一个 softmax 里竞争时，概率和边界是否仍合理。把这三件事分清后就能识破一个假成功：一个见到任何陌生照片都一律拒绝的系统，在“未知检测”上或许有点样子，但这绝不能证明它已经学会识别新增物种——拒绝只是退出，不是认识。

\subsection{多任务共享收益的度量}
\label{sec:ket-evaluation-sharing}

回到第~\ref{chap:ket-sharing} 章那三个共同交付的任务，评价多任务共享的核心是逐项看：相对各自独立学习，每个任务到底变好了还是变坏了。用 $s_k^{\mathrm{joint}}$ 和 $s_k^{\mathrm{ind}}$ 分别表示共同学习和独立学习在任务 $k$ 同一指标上的分数。为了让“正即获益”在准确率和误差两类指标上统一，用方向系数 $d_k=1$ 表示增大型指标（越大越好）、$d_k=-1$ 表示减小型指标（越小越好）；注意这里的 $d_k$ 只是方向开关，和 15.1.1 里“配对运行的差值”是两回事。要把不同单位的任务收益汇成一个数，还得先约定一个正的尺度 $c_k$ 和任务权重 $w_k$，把各项的原始差值转换成可加的标准化收益：
\begin{equation}
  g_k=d_k\frac{s_k^{\mathrm{joint}}-s_k^{\mathrm{ind}}}{c_k},
  \qquad \bar g=\sum_{k=1}^{K}w_kg_k,\quad
  w_k\geq0,\quad\sum_kw_k=1.
  \label{eq:ket-evaluation-multitask-gain}
\end{equation}
乘了方向系数后，$g_k>0$ 恒表示任务 $k$ 获益、$g_k<0$ 表示受损，方向统一了。但要警惕一个想当然：除以 $c_k$ 并不会自动造出一个有意义的“共同效用”——$c_k$ 怎么选直接决定汇总结果，它必须有明确的现实含义。合理的选法是把 $c_k$ 定成政策或效用尺度，比如“任务 $k$ 距离最低服务要求还允许变动的幅度”，这样一个单位的 $g_k$ 才代表“用掉了百分之多少的允许余量”，跨任务才可比。若随手拿独立基线分数的绝对值当 $c_k$，得到的只是“相对变化百分比”这种指数，而且当某任务基线接近零时，除以它会极不稳定。权重和尺度都必须在看到结果之前固定（否则就是按结果倒推），并且无论怎么汇总，原始单位的逐项结果都要保留备查。

用一组数字看清“平均会骗人”。设独立方案在物种准确率、病害召回率、定位平均误差上分别为 80\%、70\%、10 像素，共享方案为 84\%、68\%、9 像素。逐项看，共享带来的是：物种 $+4$ 个百分点（好）、病害 $-2$ 个百分点（坏）、定位误差减少 1 像素（好）——有一项明确退步了。若拿独立基线绝对值当尺度算相对变化指数，三项是 $5\%,-2.86\%,10\%$，等权平均约 $4.05\%$，是个漂亮的正数。但这个正数毫无业务意义：换一种指标表示或换个 $c_k$，平均值就变，它既可能掩盖病害的退步，本身也不是跨任务的共同效用，只能作为一个附条件的描述指数。

所以除了平均值，还必须报告能暴露短板的量：这组指数的最差值是 $-2.86\%$，按点估计有 $1/3$ 的任务受损。这两个量专门负责戳破“正平均”的假象，明确指出病害任务在付代价。如果病害恰是主任务，或者共同交付的合同根本不允许它退步，那么物种和定位的获益再大，也不能抵消“病害违反要求”这一事实——平均的改善在这里不构成辩护。

报告最差值时有个可比性前提：你不能在“80\%”和“10 像素”之间直接取最小，单位不同。要么都转成上面的归一化收益再取最差，要么干脆各自报告“距离最低服务要求还有多少余量”。“受损任务占比”也依赖你采用的判据，至少有两种宽严不同的口径：宽松的是按点估计 $g_k<0$ 就计为受损；严格的是只有当收益的置信区间整个落在零以下，才算“有统计证据受损”。两者描述的不是一回事，后者还要求把那些区间跨过零、方向尚不确定的任务单独列出来——既不能算获益，也不能算受损，更不能为了让成绩单好看而把它们一律记成获益。

如果价值判断上根本不接受“用一项任务的收益去补偿另一项的损失”，那就不该做加权平均，而应逐目标比较，这就引出\term{帕累托支配}（Pareto dominance）：方案 A 支配方案 B，当且仅当 A 在统一了方向的所有目标上都不差于 B，并且至少有一个目标严格更好。按这个标准，上面那个共享方案虽然物种、定位更好，但病害召回率下降了，所以它并不支配独立方案——没有任何一方在每个维度都不差。而另一个方案 $(84\%,71\%,9$ 像素) 相对 $(84\%,68\%,9$ 像素)，则在物种、定位持平的同时病害更好，按点估计确实支配后者。不过要注意比较范围：这里只看了三项任务指标，一旦把计算量、存储、延迟这些目标也加进来，就得重新逐项判断，一个在准确率上支配对手的方案，可能因算力更高而不再支配。

现实验收有时会允许每项任务有一个事先约定的小幅退步（换取其他方面的好处），但这种“容忍边界”规则和帕累托支配是两种不同的取舍哲学，不能混用。在有限样本下，严谨的做法是报告每个目标的配对差以及一组区间，再看这些区间整体落在零的哪一侧、或是否越过容忍边界。关键词是“同时覆盖”：它针对的是整组任务——当你反复按同一协议取样并为每个任务各构造一个 95\% 区间时，需要控制的是“所有这些区间一起盖住各自真实收益”的概率，而不是各管各的。每个区间单独 95\%，并不能保证整组同时 95\%（任务越多，至少一个失手的概率越大），具体的多重校正方法见第~\ref{sec:ket-evaluation-repeats}节。由此得到两个谨慎结论：点估计恰好落在容忍边界内，并不能排除真实退步更大；区间跨过边界时，结论只能是“尚不确定”。最后提醒：第14章训练中那个“共同下降方向”只保证当时各任务损失下降，它替代不了这里对最终任务指标的正式检验。

如果交付的不是单个模型，而是像 PHN 那样“按偏好在运行时生成的一整个模型族”，评价对象就必须从“一个点”扩大到“这组模型覆盖的整条取舍前沿”。偏好的集合及其抽取规则要在测试前固定，然后对每个偏好逐项报告任务成绩、有多少偏好没能满足最低要求、以及生成对应模型的费用。只展示测试后从一大堆模型里挑出的最有利的那一个，完全不能证明整个模型族在你关心的偏好范围内都具有所需取舍——那是以偏概全。

多任务比较时，预算这个参照也必须逐项钉死，否则“共享更省”很容易是错觉。设三个独立模型各用 10 GPU 小时、共 30；共享方案让各任务获得规定次数的更新，实际总共用了 24 GPU 小时。在这种设定下，你可以声称“在相同的任务更新暴露下，共享用 24 比独立用 30 省”，但绝不能说“共享只用了 10 小时”——它明明用了 24。反过来，若问题是“给定整组任务总预算 10 GPU 小时，谁效果好”，那就得让独立方案也按预定规则把这 10 小时在三个模型间分配、共享方案同样只能花 10，这才叫同总资源比较。比较时要列出的细节包括：各任务被更新的次数、标签齐全程度、缺失标签怎么处理、共享骨干的费用、调参投入。尤其记住“更新机会相同”不等于“计算相同”——不同结构走一步的代价可能差很多。

\subsection{来源能力保持与组合泛化的度量}
\label{sec:ket-evaluation-composition}

评价模型融合时，开篇那个“分别会，并不等于合起来会”的问题成为核心。把物种模型和病害模型合并后，即使它们各自的原问题仍答得很好，也不能保证系统能完成“找出指定物种里患病的花”这种联合请求。所以融合评价要拆成三个不同的主张分别测：一是原能力保持（各来源模型本来会的，融合后还会不会）；二是联合请求完成率（需要同时用到多个来源的新问题能不能答）；三是未见组合上的正确行为（训练里没出现过的属性组合能不能处理）。为衡量第一项，固定来源 $k$ 原来的评价集，用 $s_k^{\mathrm{source}}$ 记它融合前成绩、$s_k^{\mathrm{merge}}$ 记融合后系统在同一任务上的成绩，沿用上一节的方向系数，逐来源变化写成 $d_k(s_k^{\mathrm{merge}}-s_k^{\mathrm{source}})$。这里的对照必须是“该来源自己融合前”，而且除了报告来源等权平均，还必须单独列出“受损最严重的那个来源”——平均可能把某个来源的严重退化藏起来。

一个数字例子能说明“平均口径不同、结论甚至相反”。某来源在物种、病害上融合前为 90\%、86\%，融合后为 88\%、87\%：逐项看是物种 $-2$、病害 $+1$，简单等权平均由 88\% 微降到 87.5\%（似乎略变差）。但如果真实流量里 90\% 的请求其实是病害任务，按流量加权后成绩却从 $0.1\times90+0.9\times86=86.4\%$ 升到 $0.1\times88+0.9\times87=87.1\%$（反而变好）。两个平均都对，但权重含义完全不同：等权平均把两个任务看得一样重，流量平均按实际使用频率衡量。无论哪个平均上升，都不能取代“物种任务退化了 2 个点”这条逐项事实。此外还要分列两个会改变难度的条件：系统是否被告知任务身份（从而能直接选对应头）、是否被要求在统一标签空间里预测——把不同设定的分数混在一起比较没有意义。

要证明融合带来的好处超过了“只是恰当地选用了某个来源”，需要针对不同层级的操作设置直接对照，逐级加码：让最佳的单一来源回答全部请求（单来源参照，看是否真需要组合）；用一个路由器每次只选一个来源回答（检验“选得准不准”，但不产生新能力）；做预测集成，让多个来源投票或平均（多个来源都参与，但集成未必能算出需要联合推理的新关系）；最后才是真正串接模块、要求下游正确利用上游的中间结果（这才可能产生组合能力）。过程数据如来源调用日志、路由命中率有助于追踪系统行为，但要有个清醒认识：“一次请求调用了两个模型”这件事本身不等于组合成功——调了两个模型却各答各的，仍可能答不出联合问题。

怎么检验“真正的组合泛化”而不是记忆？花卉例子可用“属性组合留出”来构造：在训练里让红色、黄色这两种颜色，和斑点、条纹这两种纹理，分别在大量对象和背景上单独出现，但刻意把完整的“红色且有斑点”这一组合，从来源训练、路由器训练、提示示例直到组合验证集里全部排除。测试时要求系统找出同时满足“红色”和“斑点”两个属性的对象。关键的防作弊设计是：同一场景里还要放入“红色但条纹”和“黄色但斑点”的干扰对象——这样一来，只凭单一属性（只看红、或只看斑点）的任何取巧办法都无法答对，必须真正做合取。还要诚实界定“未见”的范围：如果底座模型在大规模预训练里早就见过红斑点的花，那你只能声称“相对本次任务训练数据留出”，不能夸大成“相对模型全部既有经验都未见过”。

组合留出还要保证“合取操作”本身在训练里有依据，否则你考的就是一个系统从没学过的推理方式、失败也不奇怪。合理的做法是：训练中提供其他属性对的合取样例（比如“黄色且条纹”“红色且条纹”），让系统有机会先学会“且”这种组合规则，再看它能否把规则迁移到没见过的“红色且斑点”上；如果训练里只有孤立的颜色标签、从无任何合取例子，却要求它无提示地执行任意逻辑组合，那属于另一个尚未建立的能力问题，不能拿来考核组合泛化。评价时应把“记忆完整模式、只用单属性、路由选择、真正的联合模块”这几类方案放在同一个留出测试集上比较，再逐步加入其他属性关系或改变指令词序，检查一次成功是否只是依赖某个固定句式。记住判断标准：模块能连上、参数能相加，只说明工程上可实现；只有在留出组合上答对，才说明真正测到了组合能力。

\subsection{全过程资源消耗的度量}
\label{sec:ket-evaluation-resources}

“适应更省钱”只是故事的一小段，紧接着要问：为了取得并长期维持这些成绩，从头到尾到底花了多少资源？很多迁移/共享方案在单步适应上便宜，却把成本转移到了别处或未来。先按时间把开销分成三类，避免漏算：来源训练和准备通常只发生一次（但可能极贵）；来源选择和目标适应是每接入一个任务就发生一次；而线上服务、后续更新和故障恢复是随时间反复发生的长期开销。沿用上一章式~\eqref{eq:ket-reuse-amortization} 的记账口径：$C_{\mathrm{prep}}$ 计入核算范围内的来源训练与准备成本，$C_{\mathrm{sel},k}$ 要计入所有候选的评分以及那些被淘汰掉的试训（失败尝试也是成本），$C_{\mathrm{adapt},k}$ 只计此前没有入过账的那部分最终适应费用（避免重复计算）。

记实际服务 $m$ 个任务、观察区间为 $[0,H]$。这里有个多任务特有的“不能重复算”的点：共享骨干对同一张图一次前向就同时给出了类别、病害、定位三项输出，这在账上只算一个推理事件，而不是三个。所以推理费应当按实际发生的推理事件集合 $\mathcal E_{\mathrm{infer}}(H)$、用每个事件的增量成本 $c_e$ 来计，不能给三个任务各记一遍前向。再计入任务 $k$ 在区间内的更新费用 $C_{\mathrm{update},k}(H)$，以及共享维护与故障恢复，得到完整的计算账目：
\begin{equation}
  \begin{aligned}
    C_{\mathrm{total}}(m,H)
    ={}&C_{\mathrm{prep}}+C_{\mathrm{shared}}(H)\\
       &+\sum_{k=1}^{m}\bigl(
        C_{\mathrm{sel},k}+C_{\mathrm{adapt},k}
        +C_{\mathrm{update},k}(H)\\
       &\hspace{24mm}+C_{\mathrm{recover},k}(H)\bigr)
        +\sum_{e\in\mathcal E_{\mathrm{infer}}(H)}c_e.
  \end{aligned}
  \label{eq:ket-evaluation-ledger}
\end{equation}
式中 $C_{\mathrm{shared}}$ 只放那些没法归到单个任务的一次性或周期性维护成本，$C_{\mathrm{recover},k}$ 包含失败重试和故障恢复的费用。记账要有“不重复”原则：同一次更新的费用一旦计入某个任务项，就不能再进共享项。有时出于内部核算需要把一次共享推理的费用分摊给各任务，这可以另定一组非负份额 $\rho_{e,k}$，对每个事件要求 $\sum_k\rho_{e,k}=1$（把这一次前向的费用在受益任务间分完），于是各任务分摊到的 $\sum_e\rho_{e,k}c_e$ 之和恰好等于真实推理总费用，不多不少。要强调分摊只改变“这笔钱算在谁头上”的归因，并不会凭空增加或减少实际发生的推理事件数。

为让账目可审计，每一笔费用都应附带：事件或版本标识、发生时间、所用设备、用途、服务了哪个任务、以及在什么条件下可被复用。一个重要原则是“成本与成败解耦”：即使某任务最终没有获益，为它花掉的选择、试训费用也照样保留在账上——失败不退款。还有个容易重复或漏算的情形：如果在来源选择阶段训出的检查点被保留下来、接着在它上面做最终适应，那适应项只加后续增量计算；如果这个检查点被丢弃、最终是重新训练的，那第二次训练就是真实发生的新成本，必须登记。只有保留这些明细，总成本才能被逐项复查而不是一笔糊涂账。

成本的核算边界有两种合理但不可混用的口径：一是只算本项目新增投入（来源模型是现成的），二是按包含来源训练在内的全生命周期核算。做方案比较时，整张对比表必须采用同一种口径——不能对“复用方案”免掉它那笔历史训练费、却让“从头训练”全额承担，那样复用天然显得便宜。准备成本应当按实际服务的任务数 $m$ 和服务期限 $H$ 摊销到每次使用上，这也是为什么复用方案在任务多时才划算；至于“将来可能还会有很多任务用它”，只能作为情景假设分析，不能拿假想的未来任务把当前账面做平。其余容易遗漏的资产——教师模型、回放记忆、数据生成器、特征缓存、检索索引、版本副本——也都要登记；并注意去重：若教师模型的某次调用已经计入训练费，就不能在推理或蒸馏项下再算同一次计算。

不同种类的资源之间往往没有公允的换算率，强行折成“总成本”会引入主观假设，所以更稳妥的是把它们分开报告。标注和审核记人工小时；计算记设备型号、运行时长和必要的运算量（如 FLOPs）；内存记峰值占用；持久存储记模型权重、样本、缓存的字节数以及保留时长。这里要区分两个容易混的存储量：随时间累积的“字节·时”（占用多久）和某一时刻的“峰值字节”（最多占多少），它们回答不同问题。故障恢复的等待、服务停顿则应另记墙钟时间（用户实际等了多久）。最后再次破除参数迷信：模型参数量少只可能意味着权重文件小一点，它既不能代替推理延迟、吞吐的实测，也不能代表总计算量小。

训练依据的制作与校验由第~\ref{part:efficient-knowledge-use}部分展开；具体服务与维护费用则按实际执行日志核算。本章用已经定义的账目、单位与去重规则比较投入，不因跨任务汇总而强行合并不同资源。

\section{评价数据与基准的选择}
\label{sec:ket-evaluation-data}

指标定清楚之后，下一个常见的致命问题是：用的数据根本测不了你声称的那件事。比如主张是“模型能推广到新地区”，可你只是把同一地区的照片随机拆成训练和测试——即便准确率很高，测试时见到的还是同一个地方的相机、背景和植物，新地区这件事压根没有被检验到。再比如，按类别人为分阶段训练可以控制学习次序，但那不等于真实世界随季节变化的数据流，测不出时间漂移下的表现。这里有个根本原则：指标只规定了“测什么量”，而评价样本必须让“主张所依赖的那个条件变化”真实出现，主张才有被证实或证伪的机会。现成基准提供了可重复的任务和数据划分，它的分数只支持“在这些特定条件下有效”的判断，绝不支持“在任意环境里都有效”。

\subsection{评价主张与数据覆盖的匹配}
\label{sec:ket-evaluation-coverage}

先看“跨领域但类别不变”这类主张怎么测。要检验花卉识别跨地区的变化，可以保持“还是分这些类”不变，但把整个采集领域留出来。DomainBed 就是为此设计的统一框架：它把多领域分类数据集、各算法和模型选择流程放进同一套评测，领域之间按照片风格、图像来源或采集环境区分，在同一套数据内变化主要落在领域、而不是类别目标。它特别区分了三种模型选择协议，权限差异必须讲清：用训练领域的数据验证、留一领域交叉验证、以及用测试领域的标签做验证。最后一种实际上动用了目标领域标签来挑超参，拥有强得多的“先知”权限，其成绩不能和“完全没有目标域标签的领域泛化”并列比较或相互借用结论。还要注意 DomainBed 检验的对象不仅是算法本身，也包括“你用什么规则选训练配置”；它并没有因此覆盖新类学习、事实修正或长期维护这些别的主张\scite{gulrajani2021}

DomainBed 的领域多少还偏“风格不同”，WILDS 则把更真实、更有社会后果的自然变化接进评价：医院、相机部署位置、地理位置、时间等，用真实应用任务做分布外划分。它包含两种结构上不同的情形，不能混为一谈：一种是训练和测试领域完全分开；另一种是领域有重叠、但其中各群体的比例发生了改变——后者你要关注的是少数/弱势群体的表现，不能简单当成“整领域留出”。和人为加噪声、扰动像素的合成偏移相比，WILDS 的变化来自真实采集条件或群体组成，应用目标仍由具体子任务决定，不是单纯的鲁棒性测试。报告时必须写明用的是哪个子集、领域/群体如何划分、官方指标是什么，不能把若干不同任务平均成一个含义不清的“WILDS 分数”。同样要克制外推：这些自然场景不覆盖所有未来偏移；在“跨医院分类”上成功，也完全没有提供“跨任务少样本学习”的证据\scite{koh2021wilds}

把主张换成“能适应一个没参与训练的学习任务”时，隔离的单位就要从“样本”升级为“整个任务”——只单独留几张照片是不够的。这里要分清任务和类别是两个不同的对象：如果系统做的仍是判断月季与鸢尾、只是把“地区乙的整个学习问题”留到最后，那检验的是“新领域下的同任务”；只有当你声称“能识别历史阶段从未见过的物种”时，才需要把这些类别本身从训练和选模中整体隔离出去。留出的层级必须和主张的层级对齐，少留一层，主张就没被真正检验。

为“适应未见任务”提供标准测床的是 Meta-Dataset：它从多套图像数据集里随机抽取少样本分类任务，任务之间在类别数、支持集大小、类别不平衡程度上都变化，部分数据集还利用类别层级构造更难的任务。它的设计让两种“未见”可以分开检验：任务内的未见类别、以及整体从未参与训练的样本来源（原协议把 Traffic Signs 和 MSCOCO 整个留作评价，其他来源则按类别划分）。但它的能力边界同样明确：这些任务本质上仍是图像分类，所以 Meta-Dataset 上的好成绩不能直接回答“能否适应病斑分割、语言生成，或任意一种新损失函数”——任务形式没变，变化的只是分类的内容\scite{ket-triantafillou2019metadataset}

未见来源有时不仅换了类别，还彻底换了成像方式——这正是 BSCD-FSL 关心的：它先用 miniImageNet 的基础类别把模型准备好，再到 CropDiseases（植物病害）、EuroSAT（卫星图）、ISIC2018（皮肤病变）、ChestX（胸部 X 光）这些成像模态迥异的数据集上评价少样本分类。研究控制了基础模型，在相同的抽样任务上直接对比元学习和普通微调，并统一采用 5 分类、每类 5/20/50 个支持样本，以便把“适应方法”这一个因素分离出来。这套设定保持了“少样本分类”的形式不变，同时大幅改变类别内容和成像领域。解读时要注意：不同数据集之间的差异是多维度的（模态、分辨率、标签结构……），不能压缩成一个单一且已被控制的“领域距离”标度；它测的是当下的跨域适应，更不能据此判断长期灾难性遗忘的情况\scite{ket-guo2020-cross-domain-few-shot}

评价多任务共享时，要特别保证“独立学习对照”是可信的，否则共享的收益说不清。自建花卉任务集应记录每张照片到底具备哪些任务的标签，以及标签缺失是否和地区、病害严重度相关（这种缺失模式本身会影响多任务结果，呼应第14章的选择偏差）。Taskonomy 在共同的室内场景图像上提供几十种密集标注，它固定输入像素、研究“一个来源任务的表示对目标任务有多大用”的迁移关系；但要强调，这个“表示可复用性”分数并不直接等于“联合训练带来的收益”——前者是单向、固定来源，后者允许任务互相改变共享表示。借用它的样本来评价共享时，还得固定每个任务用了哪些标签、共同对象、数据划分和独立预算。一个典型的归因陷阱：如果共享方案额外看到了“带物种标签的病害照片”而独立方案没有，那么性能提升可能来自多出来的监督，而不是共享机制，这两者必须用对照分开\scite{ket-zamir2018-taskonomy}

评价组合泛化的数据集，关键是真正留出“完整的关系”，而不只是换一批新的样本编号——否则模型靠记忆见过的句子就能蒙混。SCAN 是组合泛化的经典测床：它用一套有限的人工语法生成“导航指令到动作序列”的配对，其中“原语”是基本动作或指令成分，每个原语的含义和组合规则都是固定的。它提供三种难度递进的划分，各自检验不同的组合能力：随机划分（成分都见过、只是句子是新的）、长度外推（训练只见过短序列、测试更长的序列）、以及“某个原语只在单独指令里出现、测试却要求它出现在组合中”。在随机划分上得高分，完全不能说明能应付后两类。SCAN 的价值在于受控语法能干净地隔离“新组合”这个变量；但代价是它不包含真实视觉、开放语境和无限语法，所以它测的是符号层面的组合规则掌握，不是真实世界的组合理解\scite{ket-lake2018-scan}

COGS 在 SCAN 的方向上更接近真实语言：它以受控生成的英语句子为输入，要求输出表达句子语义关系的逻辑形式，从而能系统地检验几类具体的组合挑战——熟悉的词出现在新的语法角色上、句子成分的结构性重组、以及更深层的嵌套。相比 SCAN，它涉及更丰富的语言结构和语义角色，但本质上仍是规则生成的有限英语片段，因此并没有自动覆盖跨语言迁移或视觉—语言组合。报告 COGS 结果时要把同分布测试集和组合泛化集分开、并按挑战类型逐项报告：只给一个总均值，很可能把“词汇替换都对、但结构重组全错”这种关键差异掩盖掉\scite{ket-kim2020-cogs}

\begin{samepage}
评价模型融合还要根据“这些来源是怎么形成的”来选样本。Model soups 研究的是一个特定场景：从同一个预训练起点出发、用不同微调配置得到多个模型，比较其中最佳的单个模型、把权重直接平均（汤）、以及预测集成三者。要注意它的结论范围：这些模型都是为同一个目标任务微调出来的，所以“权重平均有效”这件事不能直接外推到“不同任务、不同结构模型的跨任务融合”。评价真正的多任务融合，仍须另行设置逐来源的原任务测试、联合请求测试和组合留出测试，并核查来源筛选和融合系数的确定是否偷偷用过这些测试数据（用了就是泄漏）。还有个成本误解要澄清：权重平均后只需部署一个模型，只说明推理时不必保存多个副本；训练那一大批候选模型、以及挑选平均配方的成本依然真实存在，并没有消失\scite{ket-wortsman2022-soups}

\end{samepage}

表~\ref{tab:ket-evaluation-claims}~用五项可多选标签整理上述评价对象，不将它们视为互斥类别。上块列出变化和必要留出，下块按相同行序列出基准、样本与模型访问权限与不能推出的结论。同一实验若检验多项主张，就要逐行满足相应条件；“须声明的样本与模型访问权限”列提示应报告的权限，不表示同名基准默认允许全部样本、描述或模型。

\begin{table*}[tbp]
  \centering
  \footnotesize
  \renewcommand{\arraystretch}{1.2}
  \begin{tabularx}{\linewidth}{@{}>{\raggedright\arraybackslash}p{24mm}
    *{2}{>{\raggedright\arraybackslash}X}@{}}
    \toprule
    \multicolumn{3}{@{}l@{}}{\textbf{(a) 评价主张、变化与留出}} \\
    \midrule
    对象标签（可多选） & 实际改变的条件 & 需要隔离的对象 \\
    \midrule
    \mbox{自然偏移} & 保持子任务目标，改变领域、采集条件或时期 &
      对应领域与独立对象组；时间主张另留整段时期 \\
    \addlinespace[2pt]
    \mbox{未见任务} & 变换类别、支持集规模或任务来源 &
      最终任务；未见类别主张另隔离类别，跨领域主张另隔离样本来源 \\
    \addlinespace[2pt]
    \mbox{多任务共享} & 改变任务搭配、标签覆盖或共享结构 &
      独立对象或场景组，并为各任务保留测试 \\
    \addlinespace[2pt]
    \mbox{组合泛化} & 已知成分进入新的关系、深度或次序 &
      完整组合及等价表达，同时保留基本成分覆盖 \\
    \addlinespace[2pt]
    \mbox{模型融合} & 选择来源模型参数、融合配方和服务请求分布 &
      各来源原任务测试、联合请求及组合留出 \\
    \bottomrule
  \end{tabularx}
  \par\medskip
  \begin{tabularx}{\linewidth}{@{}>{\raggedright\arraybackslash}p{24mm}
    >{\raggedright\arraybackslash}p{36mm}
    *{2}{>{\raggedright\arraybackslash}X}@{}}
    \toprule
    \multicolumn{4}{@{}l@{}}{\textbf{(b) 代表基准、访问权限与证据边界}} \\
    \midrule
    对象标签（可多选） & 样本或研究参照 & 须声明的样本与模型访问权限 & 不能据此推出 \\
    \midrule
    \mbox{自然偏移} & \mbox{DomainBed}；\mbox{WILDS} &
      来源标签；是否使用目标验证样本 & 对任意未来偏移或新任务均可适应 \\
    \addlinespace[2pt]
    \mbox{未见任务} & \mbox{Meta-Dataset}；\mbox{BSCD-FSL} &
      元训练、元验证及目标支持证据 & 有限分类任务覆盖任意任务分布 \\
    \addlinespace[2pt]
    \mbox{多任务共享} & 花卉任务集；\mbox{Taskonomy}样本 &
      每任务样本、标签及缺失情况 & 额外训练样本的收益全部来自共享机制 \\
    \addlinespace[2pt]
    \mbox{组合泛化} & \mbox{SCAN}；\mbox{COGS} &
      成分与规则可见，目标组合留出 & 符号组合成功涵盖任意多模态组合 \\
    \addlinespace[2pt]
    \mbox{模型融合} & \mbox{Model soups}；多来源任务测试 &
      来源模型、验证样本及任务头权限 & 模型参数可相加说明来源无损或新能力形成 \\
    \bottomrule
  \end{tabularx}
  \caption{主张与评价条件的对应表。两块按相同对象逐行连接：上块说明应改变和隔离什么，下块说明可参考的样本、须报告的权限及外推限制。标签可以同时选择，但各项条件不能互相代替。花卉任务集为教学设计，其余参照见正文引文；表中没有跨基准排名。}
  \label{tab:ket-evaluation-claims}
\end{table*}

\subsection{基准条件与信息来源的核查}
\label{sec:ket-evaluation-provenance}

“本次训练没有用测试文件”这句话，在预训练时代已经远远不够了——它回答不了“预训练模型是否早就接触过测试样本的内容甚至标签”。要真正核查数据重叠，需要把重叠的类型分细：完全相同的输入、输入加标签都相同、完整的任务规则相同、完整的组合相同、以及近重复样本（同一张图的裁剪、改写、不同分辨率）。这些类型对“未见”主张的破坏程度不同，必须分开记录：一张某物种的无标签照片，和一张带物种答案的测试照片，提供的是不同等级的信息（只是观测 vs. 直接监督）；模型在预训练里见过“红色”“斑点”这些词，也不等于见过“红色且有斑点”这条完整组合指令。把这些重叠类型逐一保留，你才能说清到底哪一项“留出”主张受到了影响，而不能只在论文里勾一个语焉不详的“无污染”复选框。

一份合格的核查报告至少应写清：来源语料或模型的说明文档、数据的时间范围、版本号、可访问的数据清单，以及去重的对象和方法。这里要破除对文件哈希的迷信：哈希相等能抓到字节级完全相同的文件，但抓不到图像近重复（同一张照片的不同裁剪、压缩）和文本改写（同义复述），这两类必须借助专门的近重复检索和人工抽检。报告中应给出判定阈值、发现的重叠数量、排除规则，以及“还可能漏掉什么”的风险声明。如果只检查了能下载到的公开子集，结论就只能表述为“在该公开范围内未发现所查的重叠”，不能外推到整个（可能私有的）未知语料。

预训练重叠之外，还必须登记“这一次具体运行”被允许访问什么，因为很多信息虽然不进梯度、却实实在在影响了结果。要写清：能否访问旧的原始训练数据、回放记忆和历史检查点；模型做预测时是否已知任务身份（已知身份可以直接选对应头，未知身份则要自己判断，两者难度完全不同，不能放在一起比）；标签在什么时刻到达（预测前还是预测后）；以及是否允许一次性处理一整批目标无标签输入（直推）。还要特别堵住一个常见漏洞：为来源筛选而看的目标域样本、写进提示里的示例、用来校准拒识阈值的样本、以及人工做失败分析时翻看的样本，都属于“开发访问”——哪怕它们一行梯度都没参与，它们也塑造了最终方案，必须如实登记。

每条 headline 结果都应能追溯到一份具体的访问条件清单，而不是一句模糊的“按标准协议”。一个合格的表述例如：“在新地区任务上，使用了 30 张带标签的来源选择照片、80 张适应照片和一批未标注的校准数据，但来源选择没有使用最终测试标签。”“标准协议”这种称谓本身替代不了清单——不同论文的“标准”可能差很多。下一节会把这些权限落实成具体的整体留出和样本共享规则。最后是一条研究诚信原则：如果核查中发现了足以影响主张的重叠，正确处理是重做独立评价或主动收窄结论，而不是加一条脚注说明、却在正文里保留原来那么强的断言。

对 PFN 这类“在合成任务上预训练”的方法，还要额外说明它的任务先验和合成数据生成规则，因为它的可迁移性完全由这个先验决定。一个容易被忽视的泄漏点：如果合成任务生成器或预处理配置是根据真实的开发任务调整出来的，那这些开发任务其实已经参与了“准备”，最终任务就必须另行隔离；评价时还要报告目标任务的特征分布、噪声水平、标签形式和支持集大小是否超出了准备时覆盖的范围。以 TabPFN 一类模型为例，成本和流程要分三段核算：一次性的离线任务生成与训练（很贵但摊销）、给定目标支持集时的条件推断、以及默认配置下对多种预处理做集成所产生的多次预测。使用已发布的检查点时，可以只报告本次新增的推理增量成本，但绝不能因此把历史的离线准备写成零成本，也不能把“内部跑了多次预处理预测的集成”对外说成“单次前向”。

\subsection{补充场景与失败样本的构造}
\label{sec:ket-evaluation-stress}

如果真实目标里包含稀有病害、严重遮挡或完全不同的成像方式，而评价集里净是清晰常见的图片，那么高分毫无意义——评价样本必须主动覆盖这些难条件。不同主张还要补不同的压力场景：迁移要看稀有和分布外条件；多任务共享要额外检查标签缺失和任务极度不平衡；组合泛化则必须包含没见过的完整关系和专门的干扰对象。补测哪些条件应当由“目标要求里还有哪些没被覆盖”来决定，而且纳入规则要在看到结果之前固定。这里有个统计身份的区别：测试后才发现的失败样本，可以收集起来做诊断、做回归测试，但它们不再能代表“某个预先定义总体的失败率”——因为它们是看到结果后挑出来的，那是探索性分析，不是预先设计的评测。

想搞清楚失败到底由哪个因素造成，最强的工具是“成对条件”对照：构造两个几乎相同、只差一个因素的评测条件。比如检验组合泛化时，保持颜色、纹理、对象数量和背景分布一致，只改变“某种属性关系有没有在训练中共同出现”；检验分布适应时，固定标签和样本数，只改变光照；检验多任务共享时，固定标签总量不变，只改变各任务标签的覆盖和缺失方式。这样两组之间的差异才能干净地归因到那一个被改变的因素。反过来，若一个对照同时改变了任务难度和标签数量，那观察到的差距就不能全部算在“共享结构”头上——变量没控制住。

成对条件不仅要控制变量，还要保证每个条件都有明确的任务语义，否则测的是个无意义的人造差异。组合指令要讲清逻辑连接词是“同时满足（AND）”还是“任一满足（OR）”，正确答案和计分规则必须在运行前就固定，不能事后解释。还要警惕合成数据的效度：如果合成的“病害斑”不符合真实成像规律和病理形态（比如颜色、边缘、位置都不对），那模型在它上面的表现只能说明“系统对我定义的这种图案如何响应”，不能直接当作田间真实病害识别的证据。

最后分清三类性质不同的材料，别让它们互相冒充：本章的教学数值只用于让读者复算定义、建立直觉，给它们加上误差棒也不会凭空变成真实实验；受控的合成实验用于在候选解释之间做因果区分；真实采集的样本才用于检验部署条件下的表现。每份材料都应注明自己属于哪一类、用来回答什么。若要正式补充评测样本，还需保存其生成规则、随机种子、采样范围、质量检查过程和完整的纳入/排除清单，保证所有被比较的方法接受完全相同的条件，并如实保留成功、失败和无法评价三类运行结果——只保留成功跑通的结果会系统性高估性能。

\section{对照实验与评价协议的设计}
\label{sec:ket-evaluation-protocol}

就算用的是完全相同的一套照片，“在预测之前读到标签”和“预测之后才拿到标签”也是两种根本不同的情形——前者泄漏、后者合规。所以样本一旦确定，还要把四件事固定下来：访问时机（标签和各种样本在什么时刻可见）、比较方案（基线和消融有哪些）、资源条件（预算口径）、以及重复实验覆盖的是哪一种随机性。还有一条研究纪律：如果在分析途中临时修改了来源选择规则或成功阈值，那么由此得到的新分析应当明确标注为“探索性”；若想据此确认一个新主张，必须重新独立安排一份评价数据，而不能在同一份数据上一边改规则一边下结论。

\subsection{数据划分与访问协议的设计}
\label{sec:ket-evaluation-splits}

要让数据真正检验预定的变化，必须按“与主张对应的对象”整体留出，这里有五个由浅入深的层级。样本层：隔离单张照片，并把同一植株的连拍、裁剪、增强都归为同一个对象组一起隔离（防止同株泄漏）。类别层：把最终要识别的新物种整体隔离。任务层：把整个未参与元训练和选择的预测要求隔离。领域层：隔离整个地区、设备或数据来源。时间层：严格只用过去的观测和反馈去预测未来，不允许未来信息回流。要特别说明这些层级各管各的、不能互相替代：两张照片文件不同，仍可能是同一植株；两个物种不同，仍可能来自同一地区（所以只隔离类别没隔离地区，测不了跨领域）；支持集和查询集照片分开了，它们仍可能都属于已经见过的任务族（所以测不了新任务）。

用一个例子看清层级叠加的效果。如果只在各地区、各物种内部随机拆照片，那么评价范围其实只是“已见类别、已见领域里的新样本”，测的是最浅层的泛化。把丙、丁两个物种整体留下，才第一次增加了“识别未见类别”的要求；但若这些照片仍来自已经见过的地区，新领域这件事还是没被测到。再把整个地区乙留出，才加上领域条件；再以 2026 年 7 月为界只许用过去预测未来，才又加上时间条件。多层级同时隔离会让可用训练样本变少，这是必须付出的代价，报告时应直接写明“我控制了哪几层、因此牺牲了多少数据”，而不是用“严格划分”这种含糊措辞代替具体说明。

图~\ref{fig:ket-evaluation-boundaries} 把一次严谨评测里样本的三种用途清晰分开，这是防止泄漏的主干。来源训练样本形成候选模型；目标验证集用于决定选哪个来源、用什么训练配置、何时停止；适应训练集用于真正拟合；而最终评价集只在整个方案锁定之后才用来计分，绝不提前露面。少样本/元学习中，“如何使用支持集”这类规则可以在开发任务上挑选，但一旦到了最终任务，就只能使用被明确批准的支持证据。有一种允许但必须标注的做法：如果预定协议让验证集并入适应训练集，那就要标明这一个集承担了两种用途、按唯一样本数核算标签量，并且不能再把它当成独立证据；最终评测仍必须在另外的外部数据上做。

图底部三个 3×3 网格是理解留出的绝佳示意：它们复用完全相同的九张照片，唯一改变的是“怎么画线隔离”。普通样本留出后，训练里仍含 A、B、C 三个物种和两个地区；把代表 C 物种的一整行留下，才提出“识别未见物种”的要求；把代表地区乙的一整列留下，才提出“在未见地区识别已见物种”的要求。这里变化的从来不是样本池本身，而是测试时必须跨越的边界。这正解释了前面那个反直觉结论：哪怕训练集和测试集没有一张重复照片，只要边界没画对（物种和地区都见过），就依然不足以支持类别外推或领域外推。

元学习的反馈结构比普通训练多一层，必须把任务内和任务间分开。沿用第~\ref{sec:ket-transfer-task-splits}节的记号：$\mathcal T_k\sim\Pi$ 是从任务分布抽出的学习问题，$S_k$ 用于支持适应，$Q_k$ 用于查询。关键身份要牢记：元训练任务的查询标签虽然在任务内是“考题”，但它们参与更新了学习规则，所以站在整个元学习的角度，它们是训练信号、不是独立的最终证据；元验证任务用于比较候选的学习规则，一旦据此选定并锁定规则，交给最终任务的就只是那条规则本身，而不是验证任务的数据和标签。到了最终任务，$S_k$ 提供被批准的适应证据，而 $Q_k$ 的标签只用于计分、不回流到任何开发环节。

再次强调：为最终任务新抽一组支持照片，并不等于留下了新类别或新领域，任务抽样仍必须满足上面对应层级的隔离要求。另一个必须明确声明的权限是查询输入的使用方式，它区分两种协议。直推式协议（transductive protocol）允许一次性看到同一任务的全部查询（无标签）输入并联合利用它们的结构（比如聚类、整体归一化）；归纳式协议（inductive protocol）则要求每个预测只依赖获准的支持集和当前这一个查询本身，即使你并行计算多个查询，它们之间也不能交换信息。两种协议下模型能访问的样本条件不同，所以不能只因为“支持标签数相同”就把一个直推方法和一个归纳方法直接放在一起比高低——它们的起跑线不同。

\begin{figure*}[htbp]
  \centering
  % Solid: fitting/update or labeled rule handoff; dashed: selection evidence.
% Dotted: scoring. Task collections never flow into one another.
\begin{tikzpicture}[
  x=1mm,y=1mm,font=\figurefont\footnotesize,text=BookInk,
  box/.style={draw=BookMuted,line width=0.7pt,minimum height=13mm,
    text width=27mm,align=center,inner sep=1.5mm},
  train/.style={-{Latex[length=2mm]},draw=BookTeal,line width=0.9pt},
  select/.style={-{Latex[length=2mm]},draw=BookAmber,line width=0.8pt,dashed},
  score/.style={-{Latex[length=2mm]},draw=BookBlue,line width=0.8pt,dotted},
  title/.style={anchor=west,inner sep=0pt,font=\figurefont\small\bfseries}
]
  \path[use as bounding box] (0,-178) rectangle (169,17);
  \node[title] at (0,13) {(a) 开发与独立最终评价};
  \node[box,fill=FigureInput!10] (source) at (16,0) {来源训练样本\\形成候选};
  \node[box,fill=FigureModel!10] (knowledge) at (60,0) {候选模型\\确定复用配置};
  \node[box,fill=FigureModel!10] (adapt) at (105,0) {目标适应\\执行已选规则};
  \node[box,fill=FigureOutput!10] (frozen) at (149,0) {锁定模型\\封存预测};
  \draw[train] (source.east) -- (knowledge.west);
  \draw[train] (knowledge.east) -- (adapt.west);
  \draw[train] (adapt.east) -- (frozen.west);

  \node[box,fill=FigureInput!10] (selection) at (60,-26)
    {目标验证样本\\来源、训练配置、停止};
  \node[box,fill=FigureInput!10] (support) at (105,-26)
    {目标训练样本\\拟合任务模型};
  \node[box,fill=FigureInput!10] (test) at (149,-26)
    {最终测试样本\\只读计分};
  \draw[select] (selection.north) -- (knowledge.south);
  \draw[train] (support.north) -- (adapt.south);
  \draw[score] (frozen.south) -- (test.north);
  \draw[draw=BookMuted,line width=0.6pt]
    (selection.south) -- (60,-38) -- (105,-38) -- (support.south);
  \node[anchor=north,inner sep=1mm] at (82.5,-38)
    {仅按预定协议共享};
  \node[align=left,anchor=west,inner sep=0pt,text width=30mm] at (0,-27)
    {最终结果不回流\\选择与适应};

  \node[title] at (0,-54) {(b) 元学习：任务层留出};
  \node[box,text width=31mm,fill=FigureInput!10] (metatrain) at (18,-68)
    {元训练任务\\$S_k$拟合，$Q_k$更新};
  \node[box,text width=31mm,fill=FigureModel!10] (candidates) at (62,-68)
    {候选规则};
  \node[box,text width=31mm,fill=FigureInput!10] (metaval) at (18,-88)
    {元验证任务\\试用候选规则};
  \node[box,text width=31mm,fill=FigureModel!10] (lockedrule) at (62,-88)
    {锁定规则};
  \node[box,text width=31mm,fill=FigureOutput!10] (metatest) at (62,-111)
    {最终测试任务\\$S_k$适应，$Q_k$计分};
  \draw[train] (metatrain.east) --
    node[midway,above,inner sep=1mm] {学习} (candidates.west);
  \draw[select] (candidates.south) --
    node[midway,right,inner sep=1mm] {选择} (lockedrule.north);
  \draw[select] (metaval.east) --
    node[midway,above,inner sep=1mm] {验证} (lockedrule.west);
  \draw[train] (lockedrule.south) --
    node[midway,right,inner sep=1mm] {交接规则} (metatest.north);
  \node[anchor=west,align=left,text width=33mm,inner sep=0pt] at (0,-111)
    {任务集合分开\\未见类别另做\\类别层留出};

  \node[title] at (91,-54) {(c) 连续学习：时间顺序};
  \node[box,text width=56mm,minimum height=9mm,fill=FigureOutput!10]
    (predict) at (128,-68) {状态$\sigma^{(t)}$：先预测并封存};
  \node[box,text width=56mm,minimum height=9mm,fill=FigureInput!10]
    (feedback) at (128,-85) {反馈到达：计分并按协议更新};
  \node[box,text width=56mm,minimum height=9mm,fill=FigureModel!10]
    (next) at (128,-102) {状态$\sigma^{(t+1)}$：仅服务未来};
  \draw[score] (predict.south) -- (feedback.north);
  \draw[train] (feedback.south) -- (next.north);
  \node[anchor=west,inner sep=0pt] at (91,-115)
    {独立运行重置；流内状态按协议保留};

  \node[title] at (0,-132) {(d) 同一样本池，不同留出层级};
  \foreach \shift/\kind/\heading in {0/1/样本留出,55/2/类别留出,110/3/领域留出} {
    \node[anchor=west,inner sep=0pt] at ({\shift+10},-140) {\heading};
    \foreach \r/\letter in {1/A,2/B,3/C} {
      \foreach \c in {1,2,3} {
        \ifnum\kind=1
          \ifnum\r=\c
            \fill[BookAmber!18]
              ({\shift+10+7*(\c-1)},{-146-7*(\r-1)})
              rectangle ({\shift+17+7*(\c-1)},{-153-7*(\r-1)});
          \fi
        \fi
        \ifnum\kind=2
          \ifnum\r=3
            \fill[BookAmber!18]
              ({\shift+10+7*(\c-1)},{-146-7*(\r-1)})
              rectangle ({\shift+17+7*(\c-1)},{-153-7*(\r-1)});
          \fi
        \fi
        \ifnum\kind=3
          \ifnum\c=3
            \fill[BookAmber!18]
              ({\shift+10+7*(\c-1)},{-146-7*(\r-1)})
              rectangle ({\shift+17+7*(\c-1)},{-153-7*(\r-1)});
          \fi
        \fi
        \draw[BookMuted,line width=0.5pt]
          ({\shift+10+7*(\c-1)},{-146-7*(\r-1)})
          rectangle ({\shift+17+7*(\c-1)},{-153-7*(\r-1)});
        \node[inner sep=0pt] at ({\shift+13.5+7*(\c-1)},{-149.5-7*(\r-1)})
          {\letter\c};
      }
    }
  }
  \node[anchor=west,inner sep=0pt] at (0,-174)
    {A、B、C为物种；第3列来自地区乙，其余来自地区甲；阴影为留出。};
\end{tikzpicture}

  \caption{样本用途、任务留出与时间访问边界。实线为拟合、协议更新或规则交接，虚线为选择依据，点线为只读计分；箭头不表示任务集相互并入。最终任务只接收锁定规则，查询标签不返回本轮开发；选择与适应共享须预先声明。下方九张教学照片保持相同，阴影分别标出样本、类别及领域留出。在线预测先封存，反馈只改变未来状态，不重算过去成绩。}
  \label{fig:ket-evaluation-boundaries}
\end{figure*}
\FloatBarrier

如果适应发生在一个持续接收请求的在线系统里，隔离要求还要延伸到状态：每一批“先预测、后更新”之间，模型参数、归一化统计量、优化器状态和缓存都要妥善管理，严格遵守已经声明的“先预测后更新”顺序（呼应第11章）。延迟到达的标签只能改变未来状态下的预测，绝不能回写、修改那些已经封存的历史预测成绩；开始处理一个独立任务时，则应恢复到预先约定的起点状态。至于连续序列、编辑请求和样本撤回这些更特殊的时间协议，由第~\ref{sec:ket-evolution-evaluation-access}节专门处理，本章的重点放在任务留出和共同训练的比较上。

\subsection{比较基线与消融实验的设计}
\label{sec:ket-evaluation-baselines}

设基线最大的意义是排除“替代解释”。比如来源初始化的模型成绩更高，你想归功于“来源有用”，但一个同样能解释这一结果的原因是：从头训练那条线压根没被好好优化（学习率没调、没收敛）。所以基线要专门针对这种解释来设：在保持网络结构和目标域样本一致的前提下，至少比较充分训练的从头模型、普通微调、以及冻结来源表示或随机表示这几档，才能把“来源知识”和“单纯多了优化/容量”分开。多任务则要和各任务独立学习、以及“同等总预算下的独立方案”比较。所有基线都必须被合理调参、训练到收敛——用一个没调好的随机模型当靶子来衬托自己的方法“显著更好”，是没有说服力的。

评价 PCGrad、MGDA 这类专用梯度协调器时，“单任务基线”并不是合适的直接对照——因为它连共享结构都一起取消了，差异里混入了结构因素，无法单独归因给协调器。正确的做法是固定骨干、任务头、训练样本、损失单位、初始化、基础优化器和停止规则，先放一个等权联合训练作为基准，再在预定预算内允许它搜索固定的损失权重或采样概率（把“调权”这条便宜路线也调到头）。候选范围、用什么验证集选择、搜索预算都要事先列出，测试集只评价锁定后的方案；逐任务求梯度、协调求解、淘汰配置的开销都记账，结果保留逐任务取舍而不是只报平均。这样做的现实理由来自 Xin 等人的发现：在他们测的语言任务和方法里，专用协调器给出的性能取舍并没有超过“把采样比例扫描好”的简单联合训练。所以充分调参的简单联合训练是一道必须先过的直接对照。当然这条负面证据有边界——它没覆盖所有协调器，也没测 Nash-MTL，不能推广成“所有协调器都没用”\scite{ket-xin2022-mto}

融合方法的基线要逐级配齐：至少要和“由验证集选出的最佳单一来源”以及“直接的预测集成”比较。这里要识别两种越权对照——按最终测试成绩逐任务挑最佳来源，相当于偷看答案，属于先知选择，必须明确标注；多头模型在“已知任务身份”下的集成优势，也不能把任务身份这个前提隐去不报。对结果的解读要诚实：如果你的融合在效果上略低于完整集成、但推理成本明显更低，这可以是一个合理的工程取舍，但效果和资源必须并列报告，不能只拿一个较弱的单模型当对照、就声称自己“全面占优”。若做的是模块组合，还应加上：单来源路由器对照、逐个移除模块的消融、以及用不含任务判别信息的输入做替换，检查每个成分是否真的对最终结果有贡献。

元学习的比较尤其要防止“起点不同却归因于方法”：必须先把预训练来源、骨干容量和目标支持集都保持一致，再比较三档——普通适应（用现成表示常规微调）、只用元学来的初始化但不做内层适应、以及完整的元学习适应。如果“普通适应”那一组根本没有享用相同的预训练，那两组的差异就不止“学习规则”一项，结论不成立。来源选择部分则可以比较默认顺序、随机选择、代理评分选择和等预算的目标试训，并把搜索开销一并报告。还要记住：在最终测试集上回头挑出的那个最佳来源，只能当作一个“先知诊断上界”，不能算作实际可部署的选择器性能（部署时测试标签是看不到的）。

iMAML 这类隐式元梯度方法引入了“适应解要解到多准”的新责任，比较时必须控制，否则差异可能来自求解程度而非方法本身。应固定式~\eqref{eq:ket-transfer-implicit-inner} 的内层目标和正则强度，并报告任务解的驻点残差、线性方程的求解残差、迭代次数以及实际内存和耗时。设想一方把内层解到接近最优、另一方只走了少数几步，那它们的差距很可能只是“谁解得更充分”，而不是方法优劣；为排除非线性干扰，还可以换成凸任务头，并设置同表示、同任务头的普通训练对照。PFN 的评价则要区分两个命题：一是“它对先验下的后验预测近似得准不准”，这可以和同一先验下能得到的解析解或高精度数值解比较；二是“它在真实新任务上有没有预测收益”，这仍要和普通目标训练在匹配条件下做迁移对照。前者准不代表后者一定好。

比较“凝练/蒸馏得到的训练样本”“随机原始子集”和“核心集”时，除了样本本身，标签从哪来也必须匹配，因为标签制度是一个隐藏变量。至少有三种不同的监督条件：纯人工硬标签；每个样本保存一份固定的教师软标签；每出现一个增强视图就重新向教师请求一次软标签。它们携带的信息量和成本都不同。比较时应固定教师模型版本及其训练数据权限、增强规则、向教师查询的次数、学生的训练和选择预算，并把凝练、教师准备、后续训练三类费用分别计入。Dey 等人 2026 年的研究提供了一个重要警示：在他们测的 ImageNet-1K 等设置里，只要给足教师软监督和训练算力，一个普通的随机子集就能接近被比较的凝练或核心集方法；而一旦换成硬标签，取舍又会改变。这要求我们把“监督来源（软标签）”和“样本构造方式（凝练/随机）”这两个因素分开比较，不能因为随机子集在软标签下追平了，就把结论扩大成“所有凝练方法都没有收益”\scite{ket-dey2026-soft-labels}

直接基线回答“整体好不好”，要回答“为什么好、是哪个部件在起作用”，就要靠部件干预。\term{消融实验}（ablation study）的做法是移除或替换完整方案中的某个部件，看结果如何变化。但消融并不天然等于“只改变一个变量”——移除一个部件往往连带改变了别的东西，这些伴随效应必须登记：移除回放记忆，可能同时减少了模型对旧训练样本的访问权限；移除共享骨干，改变的不仅是共享、还有参数量和计算；移除适应步骤，同时也减少了更新次数。为了区分“专门机制的作用”和“单纯多做了更新”，一个有用的对照是把消融省下的适应预算拿去做普通更新，看性能是否只是更新量带来的。只有把这些伴随变化讲清楚，消融结果才能支持机制层面的因果解释。

\subsection{资源预算与重复实验的设计}
\label{sec:ket-evaluation-repeats}

评价时预算要作为硬约束贯穿始终，而不是事后补记。如果整项开发只有一笔固定预算，那么来源筛选、那些失败的试训、以及最终适应，全都必须在这笔预算内支付（失败尝试也要花钱）；如果目标是达到同一效果，那就记录每个方案完整的过线过程，包括那些到预算用完仍未过线的运行。为了让比较可复现，来源准备的范围、共享成本要在多少任务上摊销、服务请求量、以及观察截止时刻都应事先固定。一个容易美化结果的做法是只报告那些“又便宜又成功”的运行；正确做法是：方案在预算内根本跑不完，就如实记录“不可执行”或“未完成”，这些失败记录本身就是评价结果的一部分。

“在相同资源下比较”这句话必须落实到可测量的具体量上，因为同名的资源单位可能名同实异：相同的标签数可能对应差很多的标注工时（有的标签难标）；相同的更新步数可能用了不同的批量大小和网络结构；相同的 GPU 小时可能来自不同型号的设备、在不同利用率下跑出来；相同的权重文件大小更是完全不约束训练时的内存峰值。务实的做法不是强求所有维度同时相等（那几乎不可能），而是分别设定约束：要么比标签数、要么比实际计算量、要么比内存，在某一个量上严格对齐，其余维度如实记录。那些没能对齐的差异要保留在结果表里，并相应限制你能做出的归因——不能把资源差异带来的好处说成方法优势。

只跑一次得到的配对差可能纯属运气，所以要重复实验、计算区间；但在算区间前必须先说清：你重复的到底是哪一层随机性？不同的随机来源对应不同的可推广范围。借用 15.1.1 节那组 $(6,6,4)$ 个百分点的配对收益，这里单独设一个只用于演示统计计算的教学协议：把任务、训练样本、支持样本、预先封存的测试照片、以及模型和训练配置全部固定，唯一变化的是每次运行的随机种子（初始化、数据顺序等）。每一对（迁移方案、对照方案）共享可匹配的随机条件，不同对之间独立抽种子。在这个协议下，统计上的独立单位是“完整的一对运行”，而不是其中的照片——照片在所有对里都是同一批，不能被当成重复的独立样本。再次声明这只是用相同数字演示新协议，并没有把原来那 100 张验证照片的结果改称为独立测试。

令 $d_r$ 为第 $r$ 对运行的准确率差，$\mu_d=\mathbb E_{\mathrm{run}}[d_r]$ 是在上述“其他全固定、只变运行随机”条件下，平均收益的真值。注意这个期望的下标：它只对运行随机性取平均，并不对新任务或新照片取期望——所以由它推出的区间只回答“换一批随机种子结果会怎样”，不回答“换任务、换照片会怎样”。为了在很少的运行数下也能演示区间计算，我们额外（且仅在演示中）假设这些运行差独立同分布且服从正态。前面式~\eqref{eq:ket-evaluation-run-summary} 的 $s_d$ 描述单次配对差的波动，而均值的标准误是 $\widehat{\operatorname{SE}}(\bar d)=s_d/\sqrt R$——运行次数 $R$ 越多，对平均收益的估计越精确，这正是重复实验能缩小不确定性的原因。

正态假设下，均值误差经估计标准误标准化后服从自由度$R-1$的Student $t$分布。预定95\%覆盖率，将两侧各留2.5\%，以$t_{\nu,p}$表示自由度$\nu$的$p$分位数，得到：
\begin{equation}
  \bar d\ \pm\ t_{R-1,0.975}\,
    \frac{s_d}{\sqrt R}.
  \label{eq:ket-evaluation-paired-interval}
\end{equation}
代入三次教学差值时，$R=3$、自由度为2，计算为：
\[
  \begin{aligned}
    \bar d&=\frac{16}{3},\qquad s_d=\sqrt{\frac43},\\
    \widehat{\operatorname{SE}}(\bar d)&=\frac{\sqrt{4/3}}{\sqrt3}
      =\frac23,\\
    \bar d\pm t_{2,0.975}\frac23
      &\approx5.333\pm4.303\times0.667\\
      &\approx[2.46,\,8.20]\text{个百分点}.
  \end{aligned}
\]
算完数字，最重要的是正确理解这个 95\%：它指的是“区间构造规则的长期覆盖率”——在假设成立的前提下，如果你反复抽取成组运行、每次都按这个公式造一个区间，那么长期看大约 95\% 的区间会盖住那个固定的 $\mu_d$。它绝不是说“这一次算出来的区间有 95\% 概率包含真实值”，更不是说“方案有 95\% 的概率获益”（频率派概率不这么用）。还要认清这个教学例子的局限：三个手工数字既不能用来检验正态和独立同分布假设是否成立，也不是“建议你只跑三次”。而且区间只覆盖了运行随机性，明确不包含换照片、换任务、重抽支持集、重新调参带来的不确定性——那些要靠别的实验设计覆盖。严谨的基准研究应当把这些不同的随机来源分别登记。另外，均值差和“两方案谁赢的概率”是两个不同的统计量，不能借用后者的 bootstrap 公式来给前者算区间\scite{ket-bouthillier2021-variance}

前面只重复了“同一批任务”的运行随机性。可如果你手上有多个任务，重复实验里就混着两种本质不同的变异，必须拆开：一种是同一个任务重跑多次、因种子不同产生的“运行内波动”；另一种是更换任务、因任务本身难易不同产生的“任务间差异”。以 $g_{k,r}$ 表示任务 $k$ 的第 $r$ 次配对收益，设每个任务重复 $R$ 次、共 $K$ 个等权任务。先在每个任务内部求平均 $\bar g_k$，再对任务求等权平均 $\bar g$；用两类离差平方分别刻画这两种变异：
\begin{equation}
  \begin{aligned}
    \bar g_k&=\frac1R\sum_{r=1}^{R}g_{k,r},&
    \bar g&=\frac1K\sum_{k=1}^{K}\bar g_k,\\
    s_{\mathrm{run}}^2
      &=\frac1K\sum_{k=1}^{K}
        \frac{\sum_{r=1}^{R}(g_{k,r}-\bar g_k)^2}{R-1},\\
    s_{\mathrm{task}}^2
      &=\frac{\sum_{k=1}^{K}(\bar g_k-\bar g)^2}{K-1}.
  \end{aligned}
  \label{eq:ket-evaluation-two-variations}
\end{equation}
这里要求 $R,K>1$，各自的分母对应相应的样本方差。$s_{\mathrm{run}}^2$ 平均各任务内部的运行波动（种子能解释的不稳定），$s_{\mathrm{task}}^2$ 描述不同任务均值之间的离散程度（任务选型带来的差异）。对后者要谨慎：由于每个任务只重复了有限次，任务均值本身带有估计误差，所以 $s_{\mathrm{task}}^2$ 里混着这部分估计噪声，不能直接当成“任务总体方差”的纯净估计。最要紧的是别误用这两个量：它们只是描述变异来源的统计量，并不是总均值 $\bar g$ 的标准误，不能拿它们直接构造均值的置信区间（那需要层级模型或对应的标准误推导）。

用数字区分这两种变异：两个任务分别记录三次配对收益 $(5,6,7)$ 和 $(-3,-2,-1)$ 个百分点，任务均值为 6 和 $-2$，总体等权均值为 2。两个任务内部的方差都为 1，所以 $s_{\mathrm{run}}=1$ 个百分点（同任务重跑很稳定）；而两个任务均值 6 和 $-2$ 相差很大，离差平方和为 32，故 $s_{\mathrm{task}}=\sqrt{32}\approx5.66$ 个百分点（任务间差异主导）。这个例子说明三件事：增加每个任务的重跑次数，能减小任务均值的估计波动，却不会增加你覆盖的“独立任务数”（就两个任务，跑再多遍还是两个）；而且总均值 $+2$ 看起来是正的，却完全没有消除第二个任务平均收益为 $-2$（即受损）这个事实。

t 区间依赖“运行差正态”这个可能不成立的假设。不想要这个假设时，可以用非参数方式直接从已有的独立单位近似抽样分布，这就是\term{配对自助法}（paired bootstrap）。它的核心动作是：从已有的配对运行里有放回地重抽，并且始终让两个方案使用同一组索引（抽到第 3 对就同时取两个方案的第 3 次结果），从而保住配对对应关系，再看重抽后统计量如何波动。针对“固定任务和测试样本集、只变运行”的运行均值，具体步骤如下。
\begin{enumerate}
  \item 先保存两方案每对分数与共同随机条件，固定统计量、任务权重、95\%置信水平、重抽次数$L$及种子。例如$L=10000$只表示重抽计算一万次，不是新增一万次训练实验。
  \item 每轮从$R$对运行中有放回抽$R$个索引，两方案使用同一串索引，逐对求差再平均。若汇总多个固定任务，沿用预定权重，不让重复更多的任务取得额外权重。
  \item 跑完 $L$ 轮后，取这 $L$ 个重抽均值差的第 2.5 百分位和第 97.5 百分位，作为百分位自助置信区间，并报告原始运行数、抽样单位、分位数规则和重抽次数 $L$。两个常见误解要澄清：自助区间的覆盖率是近似的（小样本下尤其不精确）；更重要的是，如果你原始只有三对运行，那么无论把这三对重抽一万次还是十万次，信息量始终来自那三个观测——自助法造不出新的独立数据，$L$ 大只是让“对这三个点的经验分布的计算”更平滑，绝不等于做了上万次独立实验。
\end{enumerate}

在多任务情境下用自助法，最容易犯的错误是破坏“共同训练”带来的任务间依赖，必须按真实的数据生成结构来重抽。如果一个元训练种子产生的共享模型同时服务全部 $K$ 个任务，那么抽中第 $r$ 次运行时，就应当把这一次的所有任务结果 $(g_{1,r},\ldots,g_{K,r})$ 整组一起带入——因为它们共享同一个模型、同一次训练，是绑在一起的；绝不能给每个任务各抽各的种子，那样会把“来自不同模型”的结果拼成一列，虚构出原实验里不存在的独立性。只有当各任务确实是独立训练、互不共享模型时，才可以分别重抽。问题若改成“固定模型、面对新测试对象时收益如何”，那独立单位变成对象组，就应重抽独立的对象组、让两方案共用组索引并保留组内照片；这种区间覆盖的是对象抽样的不确定性，并不覆盖重新训练的不确定性。

用前面那两任务的教学收益 $(5,6,7)$ 和 $(-3,-2,-1)$ 亲手检查一次重抽动作。关键前提是：两组数字的第 $r$ 项都来自同一个元训练种子 $r$，所以每一列 $(g_{1,r},g_{2,r})$ 共同代表一个模型。设想某次自助重抽抽中的种子索引串是 $(2,2,1)$，那么第一个任务应带入它的第 2、2、1 次结果 $(6,6,5)$，第二个任务必须同步带入自己的第 2、2、1 次结果 $(-2,-2,-3)$，两行列索引严格一致。如果你图省事给两个任务各抽一串不同的种子，就会把“任务1由模型2产生的结果”和“任务2由模型3产生的结果”配成一列——这在原实验里从未一起发生过，共同训练造成的相关结构就此被抹掉，区间也是错的。

如果评价目标从“固定的这 $K$ 个任务”升级为“任务分布 $\Pi$ 上的新任务”，那就需要对任务本身重抽，而这要求任务抽样确实有依据。若任务真是从 $\Pi$ 独立同分布抽取、且只评价一对固定的元模型（不重新元训练），可以重抽 $K$ 个任务索引，把每个任务完整的支持集、查询集和结果带入，重算任务平均。若还要独立重复整个元训练，就出现了任务和种子两个交叉的层级：应分别抽任务和抽种子组，且所有任务在同一次复现里共用抽中的那一串种子，以保留交叉评价的对应关系；这是在“任务抽样与种子抽样按设计相互独立”前提下的一种分组近似，不是普适公式。沿用前面的数字，若某一轮还抽中了任务串 $(2,2)$（两次都抽到第 2 个任务），那么两行都应带入第二任务的 $(-2,-2,-3)$，并按任务等权在求均值时计两次——绝不能因为任务重复出现就给它另换一串种子，那样又破坏了交叉结构。

把这个原则推广：如果重抽时连支持集或查询对象也要重新采样，就必须按实际的嵌套设计保留组内依赖（同组照片一起进出），而不能把任务、种子、照片三个层级统统摊平成一堆“独立样本”——它们本来就不独立。自助法只能从你实际抽样的那个总体里重抽：人工挑出来的两个任务，无论怎么重抽都只代表这两个任务，不会因为你反复采样就 magically 代表一个广泛的任务总体。还有一层随机性容易漏：来源/超参候选搜索过程本身的随机性，只有当你在每次复现里重做整个搜索时才被覆盖；如果固定了已经选出的配置，那搜索带来的不确定性就没有进入区间。

回到前面埋的伏笔：给每个任务各造一个 95\% 区间，并不等于“所有任务一起以 95\% 概率同时被覆盖”——任务越多，至少有一个区间失手的概率越大。如果你事先就打算对 $K$ 个任务下一个联合结论，并希望“全体同时覆盖”的总概率至少为 $1-\alpha$（取总漏覆盖概率 $\alpha=0.05$），最简单稳妥的修正是 Bonferroni 校正（Bonferroni correction）：把总预算 $\alpha$ 均分给 $K$ 个任务，每项构造 $1-\alpha/K$ 的双侧区间。具体到配对 $t$ 区间，只需把式~\eqref{eq:ket-evaluation-paired-interval} 的分位点换成 $t_{R-1,\,1-\alpha/(2K)}$。它的逻辑是联合界：即使各任务完全相关，各区间漏覆盖概率之和也至多为 $\alpha$，所以全体同时覆盖至少 $1-\alpha$——它不要求任务独立，代价是每个区间都更宽、结论更保守。

自助法也可以做同样的多重校正，把分位数改成 $\alpha/(2K)$ 和 $1-\alpha/(2K)$；但要记得自助区间本身的边际覆盖率就是近似的，做了 Bonferroni 之后得到的联合覆盖也仍是近似，无法继承 $t$ 区间那种有限样本下的较精确保证。适用范围还要讲清：这套校正针对的是“事先预定的那 $K$ 项”。如果你是在同时比较很多方法、很多预算点，或者看了数据之后才挑出要报告的任务，那么真正的比较族更大，需要相应扩大校正范围，或者更干净地另留一份独立确认样本来下结论——否则会陷入多重比较里“总能找到一个显著”的陷阱。

最后一类随机性来自“顺序”，它不能被种子重复替代。固定一个任务排列只换种子，描述的仅仅是“这个特定顺序下”的波动；要检验对不同顺序是否稳健，受控实验应明确抽取了哪些排列、是否包含对旧环境的回归测试。真实时间有个不可回避的特点：它不能倒放，你无法靠“把历史再过一遍”来增加时间维度上的独立重复——能做的是在因果边界允许的范围内重复初始化，再用不同的时间窗口或不同的实体群去检验外推是否稳健。安全性相关的攻击与删除实验沿用第3卷可信性部分的验证口径；长期维护的费用则沿用第~\ref{sec:ket-evaluation-resources} 节和式~\eqref{eq:ket-evaluation-ledger} 的资源账目，不能在汇总到“迁移/共享评价”时悄悄把它们降格成普通准确率或参数量这种轻量指标——它们各有各的严格要求。

\section{评价结果与证据边界的分析}
\label{sec:ket-evaluation-analysis}

跑完实验、拿到一堆分数，评价其实才进入最需要思考的阶段——要回头回答章首那三个疑问：到底哪些目标任务因为迁移或共享而受益、哪些受损；观察到的改善究竟由什么因素支撑（而不是假象）；以及已经测出来的结果能合理推广到什么范围。这三个问题分别由收益分布分析、机制（替代解释）分析和适用范围判断来承担，每一个都要回扣最初写下的那条主张。要破除一个观念：实验日志记录得再齐全、总分再高，也不会自动完成这些判断——数字是原料，审慎的归因和界定才是结论。

\subsection{收益分布与能力取舍的分析}
\label{sec:ket-evaluation-distribution}

做固定预算下的决策时，要先读完整条学习曲线、再看预先约定的那个终点，而不是被某一个点吸引。图~\ref{fig:ket-reuse-learning-curves} 里迁移优势从 20 标签时的 17 个百分点一路收窄到 160 标签时的 1 个点——已测各点虽仍为正，但优势大小随预算剧烈变化；第~\ref{sec:ket-reuse-prediction} 节的另一组构造更极端，曲线从早期领先直接转为后期落后，此时“只看少标签点”会得出与“看充分训练”相反的结论。因此单个早期点代表不了充分训练的结果，反之亦然。为杜绝挑点，目标任务集合和预算点都应在比较之前固定，不能因为某个预算点对自己有利就临场改用它报告。

在“达标成本”这类主张里，哪个方案更便宜还取决于一个常被忽略的变量：你实际要服务多少个任务。图~\ref{fig:ket-evaluation-utility} 沿用第~\ref{chap:ket-reuse-selection} 章的账目画清了这件事：复用方案一开始要先付一笔不菲的来源准备费（固定成本高），但此后每接入一个任务的增量费用很低；从头训练没有准备费、起步便宜，可费用随任务数线性快速增长。两条线在某个任务数处相交，交点前后“谁更便宜”发生反转——任务少时从头训练划算，任务多时复用才摊薄得回来。这说明“只用一次”和“反复用很多次”本就是两种不同的经济问题，不能放在一起评。实际分析时还要把式~\eqref{eq:ket-evaluation-ledger} 里的维护、恢复和持续服务费用补回这个简化图；若最终想合成一个同时考虑效果和计算的单一效用，所用的权重和钱—效果换算依据必须另行明确说明。

\begin{figure*}[htbp]
  \centering
  % Teaching costs from the reuse-selection chapter, equal target performance.
% C_reuse(m) = 120 + 10*m; C_scratch(m) = 30*m; all m = 1,...,10.
% Coordinate map: x = 24 + 12*m mm; y = 16 + 0.16*C mm.
\begin{tikzpicture}[
  x=1mm,y=1mm,font=\figurefont\footnotesize,text=BookInk,
  axis/.style={draw=BookInk,line width=0.7pt,-{Latex[length=1.8mm]}}
]
  \path[use as bounding box] (0,0) rectangle (169,82);
  \draw[BookTeal,line width=1pt] (29,76) -- (41,76);
  \fill[BookTeal] (35,76) circle (1mm);
  \node[anchor=west,inner sep=1mm] at (42,76) {复用};
  \draw[BookBlue,line width=1pt,dashed] (81,76) -- (93,76);
  \draw[BookBlue,fill=white,line width=0.8pt] (86,75) rectangle (88,77);
  \node[anchor=west,inner sep=1mm] at (94,76) {从头训练};

  \foreach \v in {0,100,200,300} {
    \draw[BookRule,line width=0.5pt] (24,{16+0.16*\v}) -- (144,{16+0.16*\v});
    \node[anchor=east,inner sep=1mm] at (23,{16+0.16*\v}) {$\v$};
  }
  \draw[axis] (24,16) -- (24,69);
  \draw[axis] (24,16) -- (151,16);
  \foreach \m in {0,1,2,4,6,8,10} {
    \draw[BookInk,line width=0.6pt] ({24+12*\m},16) -- ({24+12*\m},14.5);
    \node[anchor=north,inner sep=1mm] at ({24+12*\m},14.5) {$\m$};
  }
  \node[rotate=90] at (5,42) {总计算成本 $C$（GPU小时）};
  \node at (86,3) {实际服务任务数 $m$};

  \draw[BookTeal,line width=1pt]
    plot[domain=1:10,samples=10] ({24+12*\x},{16+0.16*(120+10*\x)});
  \draw[BookBlue,line width=1pt,dashed]
    plot[domain=1:10,samples=10] ({24+12*\x},{16+0.16*(30*\x)});
  \foreach \m in {1,...,10} {
    \fill[BookTeal] ({24+12*\m},{16+0.16*(120+10*\m)}) circle (1mm);
    \draw[BookBlue,fill=white,line width=0.8pt]
      ({23+12*\m},{15+0.16*(30*\m)})
      rectangle ({25+12*\m},{17+0.16*(30*\m)});
  }
  \draw[BookMuted,line width=0.7pt,dotted] (96,16) -- (96,44.8);
  \draw[BookMuted,line width=0.7pt] (96,46) -- (96,57);
  \node[anchor=south,inner sep=1mm] at (96,57) {$m=6$持平};
  \node[anchor=north west,inner sep=0pt,text=BookMuted] at (31,68)
    {成本越低越好};
\end{tikzpicture}

  \caption{同效果条件下，实际服务任务数$m$改变费用排序的教学构造。复用为$C_{\mathrm{reuse}}=120+10m$，从头训练为$C_0=30m$，单位GPU小时；相同推理费相消，维护等费用暂略。两者在$m=6$均为180，此后复用较低。点对应整数任务数，连线只引导观察，不表示额外测量。}
  \label{fig:ket-evaluation-utility}
\end{figure*}

分析融合结果时，要把“来源能力退化”和“组合留出上的失败”并列摆出来，不能只报联合请求成功率；而涉及跨时间阶段修改模型的部分，则按第~\ref{chap:ket-evolution-evaluation} 章的要求另行验证，不能在本章顺手带过。如果各目标之间确实出现了取舍（一个变好、另一个变差），判断标准不应是事后看总分，而要回到第~\ref{chap:ket-sharing} 章事先约定好的主辅关系或共同交付要求：哪些任务是必须保住的、哪些允许小幅退步。报告时把对象分成三类列出——受益的、受损的、以及证据方向尚不确定的。一个正的总体收益数字，并不会改变这些事先约定的硬约束是否被满足。

\subsection{机制贡献与替代解释的分析}
\label{sec:ket-evaluation-mechanisms}

看到“共享方案比独立方案准”时，机制分析的第一反应应当是：这个改善真的来自共享吗？还是来自任何一个混进来的额外优势？如果共享方案多看了标注照片，那么“额外监督”就是一个和“共享机制”竞争的替代解释，当前比较无法把改善全算在共享头上。类似的可能解释还有：它用了更强的预训练、花了更多训练计算、融合后网络容量变大了——任何一项都足以单独带来提升。要把机制讲清楚，就得把完整方法、足够强的基线和关键消融放进“同一张样本访问与预算对照表”里，逐项检查哪些条件确实被固定了、哪些差异其实还没被分离开；只有排除了这些替代解释，结论才能写成“共享机制导致了改善”，否则只能写“这一整套方案表现更好”。

来源选择用的“代理评分”便宜，但必须拿它和真实的适应收益对照，不能默认便宜指标就准。可以检验几件具体的事：固定适应预算后，代理评分高的来源是否真的领先；做一次真实试训会不会改变来源排名；那些被代理早早淘汰的来源，会不会只是早期收敛慢、其实终点更好。还要认清相关性的边界：代理排序和最终真实排序相关，只说明“候选之间的相对次序大体一致”，它完全不提供绝对水平的保证——可能所有候选按代理排得明明白白，却统统不如干脆不迁移的对照。另外，靠完整试训拿到更高的终点成绩，往往也意味着花了更多标签和算力，这要按 15.1.1 节的口径分别评价“固定预算下的收益”和“达标所需成本”，不能只报终点高。

观察性结论最容易被误读成因果。比如你发现“梯度越一致的任务对，共享收益越高”，这如果只是从已有实验里观察到的相关，就不能反过来说“提高梯度一致性就能提高收益”——很可能是第三个因素（任务本来就相似）同时导致了两者。要走向因果，得做干预：固定其他条件，主动改变任务搭配，直接检验“操纵一致性是否带来收益”这个更具体的假说；并且要在新地区的真实工作负载上复现，才增加部署条件下的证据。对证据强度也要有统一标准：即便是近期的预印本，同样需要公开配置、可重复、做消融、并跨不同数据来源复核；发表时间新、框架名字响亮，都不能替代这些检查。如果结论只在少数几个模型上验证过，那就老老实实地把结论范围限定在这些模型内，不外推。

\subsection{适用范围与证据缺口的分析}
\label{sec:ket-evaluation-limits}

结论的适用范围应当像填写条件卡片一样精确，而不是一句笼统的“迁移有效”。一条严谨的少样本结论应写成：“在已测的这类类别配置的图像分类任务上，当提供类别描述、每任务给五个支持样本、并采用与对照相同的预训练起点时，该方法达到目标水平所需的标签比普通适应更少。”这句话里的每个限定词都在承担边界：支持样本存在（不是无监督）、任务是图像分类（不是任意任务）、起点和预算都已固定。删掉任何一个限定，把它缩写成“实现了无需监督的迁移”，都是过度声称。还要回到第13章提过的观测限制：即便目标输入分布看起来相同，背后的标签规则也可能不同；所以只靠无标签代理评分选来源，并没有直接确认目标任务上的真实收益，增加再多同类无标签评分也无法单独判定哪个来源的目标风险更低，这一点的理论依据可参考\BookChapterRef{chap:transfer-learning-theory}{卷3《基础、理论与可信性》的“迁移学习的可学习性”章}。

评价得到的证据有强有弱，至少要区分三档，别把它们当成同等强度：一是“在测试中没有检出失败/负迁移”，这是最弱的一档，它可能只是因为测的任务太少或检验功效不足，没发现不等于不存在；二是“在指定测试集上达到了标准”，这只对那个测试分布负责；三是“在明确列出的假设下得到了理论保证”，这最强但结论严格依赖假设成立。还要注意统计显著性和实际价值是两回事：未检出收益，既可能因为真的差异很小，也可能因为任务或运行次数不足、检验不够灵敏；反过来说，一个统计上能检出的极小差异，也可能小到对实际使用毫无意义。所以报告要把差异幅度、不确定性（区间）、成功要求和尚未测试的范围并列写出，不能只丢一个“显著/不显著”了事。

当证据不足以支持原来的强主张时，正确而专业的反应有两种：要么主动把结论收窄到确实被支持的条件内，要么针对具体缺口安排新的评价。缺口的类型决定补什么实验：领域覆盖不够，就做新的领域留出；费用有遗漏，就补齐全过程账目；组合任务失败，就设计成对的组合条件并明确规则。要特别警惕一种“拼装式高分”：把不同访问权限下测的结果、不同任务上的成绩、不同分母算出的比率拼成一个漂亮数字——这种高分无法补足其中任何一项的证据缺陷，三项基础主张（迁移改善目标学习、关键共享任务没有受损、组合能力超出单纯来源选择）只要有一项证据不足，整体结论就不能成立。全章最终回到开篇的三个疑问，要求分别给出判断，并附上取得和长期维持这些结果所需的总成本，这才构成一个完整而诚实的评价闭环。

\section{本章小结}
\label{sec:ket-evaluation-summary}

全章可以浓缩成几条判断原则。判断迁移是否有益，只有两条干净的途径：在相同目标证据和相同预算下比较目标成绩，或比较各方达到同一要求所付出的全部代价。判断共享是否有益，则必须保留逐任务的收益、受损和不确定性，绝不能用一个平均分数代替。而来源能力保持、联合请求完成、未见组合泛化是三个彼此独立的主张，必须分别验证——它们都不能从“模块在工程上可以连接”或“模型参数可以相加”这种可行性事实直接推出来：连得上不等于答得对。

数据和协议层面，任务、类别、领域、完整组合这四种留出分别检验不同范围的泛化，必须与主张对齐、不可互相替代；支持样本可以参与目标适应，但最终查询标签绝不回流到开发环节。对照层面，独立学习、充分调参的简单联合训练、最佳单来源和预测集成分处不同层级，各自排除一类替代解释，缺一不可。最终能得出什么可用结论，是由指标、全过程费用账目、配对重复实验和明确的结论边界共同决定的。至于模型上线后的持续维护、能力保持、遗忘与撤回这些随时间产生的特殊责任，则不在本章闭环之内，由第~\ref{chap:ket-evolution-evaluation} 章另行核对。
